\documentclass[10pt,a4paper]{amsart}
\usepackage{amsmath,amssymb,amsthm,graphicx}
\usepackage{verbatim}
\newenvironment{reference}{\comment}{\endcomment}
%\newenvironment{reference}{}{}
\newenvironment{slogan}{\comment}{\endcomment}
\newenvironment{history}{\comment}{\endcomment}

% For commutative diagrams we use Xy-pic
\usepackage[all]{xy}

% We use 2cell for 2-commutative diagrams.
\xyoption{2cell}
% Theorem environments.
%
\theoremstyle{plain}
\newtheorem{theorem}[subsection]{Theorem}
\newtheorem{proposition}[subsection]{Proposition}
\newtheorem{lemma}[subsection]{Lemma}

\theoremstyle{definition}
\newtheorem{definition}[subsection]{Definition}
\newtheorem{example}[subsection]{Example}
\newtheorem{exercise}[subsection]{Exercise}
\newtheorem{situation}[subsection]{Situation}

\theoremstyle{remark}
\newtheorem{remark}[subsection]{Remark}
\newtheorem{remarks}[subsection]{Remarks}

\numberwithin{equation}{subsection}

% Macros
%
\def\lim{\mathop{\mathrm{lim}}\nolimits}
\def\colim{\mathop{\mathrm{colim}}\nolimits}
\def\Spec{\mathop{\mathrm{Spec}}}
\def\Hom{\mathop{\mathrm{Hom}}\nolimits}
\def\Ext{\mathop{\mathrm{Ext}}\nolimits}
\def\SheafHom{\mathop{\mathcal{H}\!\mathit{om}}\nolimits}
\def\SheafExt{\mathop{\mathcal{E}\!\mathit{xt}}\nolimits}
\def\Sch{\mathit{Sch}}
\def\Mor{\mathop{\mathrm{Mor}}\nolimits}
\def\Ob{\mathop{\mathrm{Ob}}\nolimits}
\def\Sh{\mathop{\mathit{Sh}}\nolimits}
\def\NL{\mathop{N\!L}\nolimits}
\def\CH{\mathop{\mathrm{CH}}\nolimits}
\def\proetale{{pro\text{-}\acute{e}tale}}
\def\etale{{\acute{e}tale}}
\def\QCoh{\mathit{QCoh}}
\def\Ker{\mathop{\mathrm{Ker}}}
\def\Im{\mathop{\mathrm{Im}}}
\def\Coker{\mathop{\mathrm{Coker}}}
\def\Coim{\mathop{\mathrm{Coim}}}

% Boxtimes
%
\DeclareMathSymbol{\boxtimes}{\mathbin}{AMSa}{"02}

%
% Macros for moduli stacks/spaces
%
\def\QCohstack{\mathcal{QC}\!\mathit{oh}}
\def\Cohstack{\mathcal{C}\!\mathit{oh}}
\def\Spacesstack{\mathcal{S}\!\mathit{paces}}
\def\Quotfunctor{\mathrm{Quot}}
\def\Hilbfunctor{\mathrm{Hilb}}
\def\Curvesstack{\mathcal{C}\!\mathit{urves}}
\def\Polarizedstack{\mathcal{P}\!\mathit{olarized}}
\def\Complexesstack{\mathcal{C}\!\mathit{omplexes}}
% \Pic is the operator that assigns to X its picard group, usage \Pic(X)
% \Picardstack_{X/B} denotes the Picard stack of X over B
% \Picardfunctor_{X/B} denotes the Picard functor of X over B
\def\Pic{\mathop{\mathrm{Pic}}\nolimits}
\def\Picardstack{\mathcal{P}\!\mathit{ic}}
\def\Picardfunctor{\mathrm{Pic}}
\def\Deformationcategory{\mathcal{D}\!\mathit{ef}}
\usepackage{cleveref}
\crefname{originalreference}{original reference}{original references}
\Crefname{originalreference}{Original reference}{Original references}
\makeatletter
\newlabel{homology-section-truncations}{{0118}{source}{}{sourcepointer-1}{}}
\newlabel{homology-section-truncations@cref}{{[originalreference][0][]0118}{[1][0][]source}{}{}{}}
\makeatother
\crefname{originallemma}{original lemma}{original lemmas}
\Crefname{originallemma}{Original lemma}{Original lemmas}
\makeatletter
\newlabel{lemma-poincare}{{07LC}{source}{}{sourcepointer-2}{}}
\newlabel{lemma-poincare@cref}{{[originallemma][0][]07LC}{[1][0][]source}{}{}{}}
\makeatother
\crefname{originalreference}{original reference}{original references}
\Crefname{originalreference}{Original reference}{Original references}
\makeatletter
\newlabel{homology-section-filtered-complex}{{012K}{source}{}{sourcepointer-3}{}}
\newlabel{homology-section-filtered-complex@cref}{{[originalreference][0][]012K}{[1][0][]source}{}{}{}}
\makeatother
\title{Relative Poincare lemma: de Rham complex of the divided power polynomial extension}
\author{The Stacks Project Authors}
\date{stacks-project:crystalline.tex:local-snapshot-20261012; GNU FDL 1.2 or later}
\begin{document}
\maketitle
\noindent\textit{Editorial note.} All delivered mathematical text is the author's own, copied verbatim from the pinned Stacks Project chapter file crystalline.tex; no statement or proof step was rewritten, reordered or removed.\par\medskip
\noindent\textit{Editorial note.} Span recomputation by label pairing (LANE\_BRIEF 10.2): in this chapter the author's proof environment follows the statement environment as a standalone block; the spec covers the complete environment pair (statement span, proof span) as recorded in the pool candidate.\par\medskip
\ifdefined\hypertarget\hypertarget{sourcepointer-1}{}\fi
\noindent\textit{Source reference.} Original reference 0118: sources/stacks-chow/tags.txt, source lines 1354--1354. Cross-chapter standard prerequisite cited by the author inside the counted span; it is used as a pointer and its own proof is not part of this item.\par\medskip
\ifdefined\hypertarget\hypertarget{sourcepointer-2}{}\fi
\noindent\textit{Source reference.} Original lemma 07LC: sources/stacks-chow/tags.txt, source lines 9687--9687. Cross-chapter standard prerequisite cited by the author inside the counted span; it is used as a pointer and its own proof is not part of this item.\par\medskip
\ifdefined\hypertarget\hypertarget{sourcepointer-3}{}\fi
\noindent\textit{Source reference.} Original reference 012K: sources/stacks-chow/tags.txt, source lines 1402--1402. Cross-chapter standard prerequisite cited by the author inside the counted span; it is used as a pointer and its own proof is not part of this item.\par\medskip
\begin{lemma}
\label{lemma-relative-poincare}
Let $A$ be a ring. Let $(B, I, \delta)$ be a divided power ring
where $B$ is an $A$-algebra.
Let $P = B\langle x_i \rangle$ be a divided power polynomial
ring over $B$ with divided power ideal $J = IP + B\langle x_i \rangle_{+}$
as usual. Let $M$ be a $B$-module endowed with an integrable connection
$\nabla : M \to M \otimes_B \Omega^1_{B/A, \delta}$. Then the map of
de Rham complexes
$$
M \otimes_B \Omega^*_{B/A, \delta}
\longrightarrow
M \otimes_P \Omega^*_{P/A, \delta}
$$
is a quasi-isomorphism. Let $D$, resp.\ $D'$ be the $p$-adic completion of
$B$, resp.\ $P$ and let $\Omega^i_D$, resp.\ $\Omega^i_{D'}$ be the $p$-adic
completion of $\Omega^i_{B/A, \delta}$,
resp.\ $\Omega^i_{P/A, \delta}$. Let $M$ be a $p$-adically complete
$D$-module endowed with an integral connection
$\nabla : M \to M \otimes^\wedge_D \Omega^1_D$.
Then the map of de Rham complexes
$$
M \otimes^\wedge_D \Omega^*_D
\longrightarrow
M \otimes^\wedge_D \Omega^*_{D'}
$$
is a quasi-isomorphism.
\end{lemma}

\begin{proof}
Consider the decreasing filtration $F^*$ on $\Omega^*_{B/A, \delta}$
given by the subcomplexes
$F^i(\Omega^*_{B/A, \delta}) = \sigma_{\geq i}\Omega^*_{B/A, \delta}$.
See Homology, Section \ref{homology-section-truncations}.
This induces a decreasing filtration $F^*$ on $\Omega^*_{P/A, \delta}$
by setting
$$
F^i(\Omega^*_{P/A, \delta}) =
F^i(\Omega^*_{B/A, \delta}) \wedge \Omega^*_{P/A, \delta}.
$$
We have a split short exact sequence
$$
0 \to \Omega^1_{B/A, \delta} \otimes_B P \to
\Omega^1_{P/A, \delta} \to
\Omega^1_{P/B, \delta} \to 0
$$
and the last module is free on $\text{d}x_i$. It follows from this that
$F^i(\Omega^*_{P/A, \delta}) \to \Omega^*_{P/A, \delta}$ is a termwise
split injection and that
$$
\text{gr}^i_F(\Omega^*_{P/A, \delta}) =
\Omega^i_{B/A, \delta} \otimes_B \Omega^*_{P/B, \delta}
$$
as complexes. Thus we can define a filtration $F^*$ on
$M \otimes_B \Omega^*_{B/A, \delta}$ by setting
$$
F^i(M \otimes_B \Omega^*_{P/A, \delta}) =
M \otimes_B F^i(\Omega^*_{P/A, \delta})
$$
and we have
$$
\text{gr}^i_F(M \otimes_B \Omega^*_{P/A, \delta}) =
M \otimes_B \Omega^i_{B/A, \delta} \otimes_B \Omega^*_{P/B, \delta}
$$
as complexes.
By Lemma \ref{lemma-poincare} each of these complexes is
quasi-isomorphic to $M \otimes_B \Omega^i_{B/A, \delta}$ placed in degree $0$.
Hence we see that the first displayed map of the lemma is a morphism of
filtered complexes which induces a quasi-isomorphism on graded pieces. This
implies that it is a quasi-isomorphism, for example by the spectral sequence
associated to a filtered complex, see
Homology, Section \ref{homology-section-filtered-complex}.

\medskip\noindent
The proof of the second quasi-isomorphism is exactly the same.
\end{proof}

\end{document}
